\noindent\textbf{样例 1}（篡改类型：实体；\method{} 预测 $\hat y=0.00$）\par
{\small\begin{quote}
[1] \textbf{〔被篡改〕}【华为隐形指纹技术正式亮相】\quad\textbf{〔原句〕}【vivo隐形指纹技术正式亮相】\par
[2] 全面屏手机时代即将到来。\par
[3] 如此一来前后指纹已成为过去式，所有手机设计上又将站在同一起跑线了，大家觉得屏幕内指纹解锁帅不帅？秒拍视频\par
\end{quote}}
{\small\begin{tabularx}{\linewidth}{@{}clXccc@{}}
\toprule
句 & 类型 & 声明与两方论证 & $v_k$ & $g_k$ & $p_k$ \\
\midrule
-- & 整体 & \textbf{这条消息整体的可信度}\newline $+$ 全面屏手机的发展趋势确实存在。\newline $-$ 消息中提到的华为隐形指纹技术未得到官方证实。 & 0.74 & 0.23 & 0.00 \\
1 & 事实 & \textbf{华为隐形指纹技术正式亮相}\newline $+$ 华为经常发布新技术，此消息可能为谣言。\newline $-$ 华为官方未发布过此消息。 & 0.59 & 0.62 & 0.01 \\
2 & 观点 & \textbf{全面屏手机时代即将到来}\newline $+$ 手机厂商都在研发全面屏，此说法有一定合理性。\newline $-$ 全面屏技术仍在发展，尚未普及。 & 0.45 & 0.54 & 0.00 \\
\bottomrule
\end{tabularx}}
\par\medskip
\noindent\textbf{样例 2}（篡改类型：实体；\method{} 预测 $\hat y=0.00$）\par
{\small\begin{quote}
[1] 【天津健康码亮码超2100万次集成新功能后将沿用】\par
[2] 很多人关心疫情过后，天津健康码是否还会沿用？\par
[3] 在3月19日召开的新闻发布会上，天津市大数据管理中心主任边柯柯就此回应称，天津健康码将联合津心办一网通办平台，集成更多便民服务应用，继续沿用下去。\par
[4] \textbf{〔被篡改〕}边柯柯介绍说，天津将积极推动电子通行证在互联网+医疗健康相关场景的应用，实现签约预约\quad\textbf{〔原句〕}边柯柯介绍说，天津将积极推动健康码在互联网+医疗健康相关场景的应用，实现签约预约\par
\end{quote}}
{\small\begin{tabularx}{\linewidth}{@{}clXccc@{}}
\toprule
句 & 类型 & 声明与两方论证 & $v_k$ & $g_k$ & $p_k$ \\
\midrule
-- & 整体 & \textbf{这条消息整体的可信度}\newline $+$ 声明直接引用了官方发言，具有较高的可信度。\newline $-$ 消息中提及的具体日期和发言人姓名不明确，且部分内容为预测性意见，缺乏具体细节。 & 0.52 & 0.50 & 0.01 \\
3 & 引述 & \textbf{天津市大数据管理中心主任边柯柯在3月19日的新闻发布会上表示，天津健康码将集成更多便民服务应用，继续沿用下去。}\newline $+$ 该声明直接引用了官方发言，具有较高的可信度。\newline $-$ 新闻发布会的具体日期和发言人姓名未明确提及，可能存在信息不准确的风险。 & 0.23 & 0.82 & 0.01 \\
4 & 观点 & \textbf{天津将积极推动电子通行证在互联网+医疗健康相关场景的应用。}\newline $+$ 这是官方对未来发展的积极展望，符合政策导向。\newline $-$ 该说法是预测性意见，缺乏具体实施细节和时间表，可信度较低。 & 0.64 & 0.89 & 0.03 \\
\bottomrule
\end{tabularx}}
\par\medskip
